\phantomsection
\pdfbookmark[1]{Abstract}{sec:abstract}
\begin{abstract}
This report describes a UART-controlled coarse-grained reconfigurable array
with a default $4\times4$ mesh and configurable rectangular geometries up to
sixteen 16-bit processing elements. Each element has one result register;
the host configures operations, supplies edge operands and requests a count
of logical steps. We describe diagonal operations, local MAC reductions,
matrix products with stationary coefficients and host reduction, and prefix
scans. A public C API and declarative front end share the kernel engines.
Independent scalar references, emulator/RTL protocol comparison and UART
simulation verify functional behaviour across eight geometries. A reproducible
experiment compares current protocol v2 and v3 using byte and transaction
counters with fixed workloads. These counts isolate fused execution without
claiming physical latency or CPU speedup. The report also specifies enabled
capture spacing and the Vivado timing gates. A real synthesis and post-route
STA run passes at a 10 ns constraint for the default mesh; physical board
operation and end-to-end acceleration remain unmeasured.
\end{abstract}
